% Abhakseite „Das kann ich“ – erzeugt von rahmen.py aus den Titeln der Quelltexte
\clearpage
\hypertarget{abhaken}{}\einheitenkopf{Das kann ich}
{\small\setlength{\parskip}{0pt}
\abhakgruppe{Kennst du schon}
\abhak{1}{Ich kann besondere Punkte eines Graphen ablesen.}
\abhak{2}{Ich kann Gleichungen durch Ausklammern und mit der Lösungsformel lösen.}
\abhak{3}{Ich finde den Fehler beim Lösen durch Ausklammern.}
\abhak{4}{Ich kann Ableitungen bilden.}
\abhak{5}{Ich kann Funktionswerte berechnen und prüfen, ob ein Punkt auf einem Graphen liegt.}
\abhak{6}{Ich kann am Term erkennen, wie ein Graph grob verläuft.}
\abhak{7}{Ich kann an einem Graphen ablesen, wie steil er ist.}
\abhakgruppe{Einheit 1 · Monotonie und erste Ableitung}
\abhak{8}{Ich kann am Vorzeichen der Ableitung ablesen, ob ein Graph steigt oder fällt.}
\abhak{9}{Ich kann die Intervalle berechnen, in denen ein Graph steigt oder fällt.}
\abhak{10}{Ich kann am Term der Ableitung zeigen, dass ein Graph überall steigt oder überall fällt.}
\abhak{11}{Ich finde den Fehler bei den Monotonieintervallen.}
\abhak{12}{Ich kann begründen, wann ein Graph überall steigt.}
\abhak{13}{Ich kann mit der Monotonie eine Aussage über zwei Modelle beurteilen.}
\abhakgruppe{Einheit 2 · Extrempunkte}
\abhak{14}{Ich erkenne, ob eine Stelle gegeben ist oder gesucht wird.}
\abhak{15}{Ich erkenne, welcher Nachweis für die Art eines Extrempunkts passt.}
\abhak{16}{Ich kann die Stellen berechnen, an denen ein Graph eine waagerechte Tangente hat.}
\abhak{17}{Ich kann Hoch-, Tief- und Sattelpunkte berechnen.}
\abhak{18}{Ich kann nachweisen, dass an einer gegebenen Stelle ein Hoch- oder Tiefpunkt liegt.}
\abhak{19}{Ich kann mit berechneten Hoch- und Tiefpunkten weiterrechnen.}
\abhak{20}{Ich finde den Fehler bei der Art eines Extrempunkts.}
\abhak{21}{Ich kann ohne Rechnung begründen, ob an einer Stelle ein Hoch- oder Tiefpunkt liegt.}
\abhak{22}{Ich kann mit der Ableitung begründen, welche Hoch- und Tiefpunkte ein Graph hat.}
\abhak{23}{Ich kann Gipfelpunkte in einer Sachaufgabe berechnen und vergleichen.}
\abhakgruppe{Einheit 3 · Krümmung und Wendepunkte}
\abhak{24}{Ich kann die zweite und die dritte Ableitung bilden.}
\abhak{25}{Ich kann Wendepunkte berechnen.}
\abhak{26}{Ich kann Wendepunkte berechnen – weiter.}
\abhak{27}{Ich kann Geraden durch Wendepunkte untersuchen.}
\abhak{28}{Ich finde den Fehler beim Krümmungsverhalten.}
\abhak{29}{Ich kann begründen, was an einem Wendepunkt geschieht.}
\abhak{30}{Ich kann die Krümmung im Sachzusammenhang deuten.}
\abhakgruppe{Einheit 4 · Graph und Ableitungsgraph}
\abhak{31}{Ich kann am Graphen von $f$ ablesen, wo $f'$ positiv, negativ oder null ist.}
\abhak{32}{Ich kann den Graphen der Ableitung zum Graphen von $f$ zeichnen.}
\abhak{33}{Ich kann einen möglichen Graphen zu vorgegebenen Eigenschaften skizzieren.}
\abhak{34}{Ich kann aus dem Graphen einer Differenzfunktion die Lage zweier Graphen ablesen.}
\abhak{35}{Ich kann Aussagen über einen Graphen und seine Ableitung beurteilen.}
\abhak{36}{Ich finde den Fehler beim Lesen eines Ableitungsgraphen.}
\abhak{37}{Ich kann Zusammenhänge zwischen $f$ und $f'$ begründen.}
\abhak{38}{Ich kann aus dem Graphen einer Änderungsrate den Verlauf eines Bestands ablesen.}
\abhakgruppe{Einheit 5 · Kurvenuntersuchung im Sachzusammenhang}
\abhak{39}{Ich erkenne, was eine Sachfrage mathematisch bedeutet.}
\abhak{40}{Ich kann den Verlauf eines Graphen im Sachzusammenhang beschreiben.}
\abhak{41}{Ich kann größte Werte und stärkste Änderungen im Sachzusammenhang berechnen.}
\abhak{42}{Ich kann größte Werte und stärkste Änderungen im Sachzusammenhang berechnen – weiter.}
\abhak{43}{Ich finde den Fehler bei der stärksten Zunahme.}
\abhak{44}{Ich kann Ergebnisse im Sachzusammenhang begründen.}
\abhak{45}{Ich kann die Maße eines umschließenden Quaders aus einem Profil bestimmen.}
\abhak{46}{Ich kann Steigung und Neigungswinkel im Sachzusammenhang berechnen.}
\abhak{47}{Ich kann Werte am Graphen im Sachzusammenhang ablesen und deuten.}
}
